\section*{Chapter 7 추가 기호와 용어}

\renewcommand{\arraystretch}{1.2}
\begin{longtable}{>{\raggedright\arraybackslash}p{0.25\textwidth}>{\raggedright\arraybackslash}p{0.65\textwidth}}
\toprule
\textbf{기호·용어} & \textbf{뜻} \\
\midrule
\endfirsthead
\toprule
\textbf{기호·용어} & \textbf{뜻} \\
\midrule
\endhead
\bottomrule
\endfoot

$C_G(g)$ & 원소 $g$와 가환하는 $G$의 모든 원소가 이루는 중심화군. \\
$N_G(H)$ & 부분군 $H$를 켤레작용으로 보존하는 원소들이 이루는 정규화군. \\
$\operatorname{Cl}_G(g)$ & $g$의 $G$-켤레류 $\{xgx^{-1}:x\in G\}$. 그 크기는 $[G:C_G(g)]$. \\
단순군 & 비자명한 진정한 정상부분군을 갖지 않는 비자명군. \\
완전군 & 교환자부분군이 군 전체와 같은 군 $[G,G]=G$. \\
$K_n=k(C_1,\dots,C_n)$ & 일반 $n$차다항식의 대수적으로 독립인 계수들이 생성하는 유리함수체. \\
일반다항식 $P_n$ & $X^n+C_1X^{n-1}+\cdots+C_n\in K_n[X]$. 분해체의 갈루아군은 $S_n$. \\
보편 근호공식 & 일반다항식의 모든 근을 포함하는 계수체 위 근호확장이 존재한다는 체론적 조건. \\
아벨--루피니 정리 & 표수 $0$에서 차수 $n\ge5$의 일반다항식은 근호로 풀리지 않는다는 정리. \\

\end{longtable}
